\chapter{Operating System Security}

\section{Introduction}

Operating System (OS) security focuses on protecting system resources, processes, memory, files, devices, and users from unauthorized access, misuse, and attacks.

Modern operating systems such as Linux, Windows, macOS, Android, and iOS implement multiple layers of defense to provide confidentiality, integrity, and availability.

For Software Development Engineers (SDEs), OS security knowledge is important because security vulnerabilities often arise from improper privilege handling, memory misuse, weak authentication mechanisms, and insecure system design.

\section{Learning Objectives}

After completing this chapter, you should be able to:

\begin{itemize}
\item Understand core OS security concepts.
\item Explain authentication and authorization mechanisms.
\item Compare ACLs and capabilities.
\item Understand memory protection techniques.
\item Explain ASLR and DEP.
\item Understand sandboxing and secure boot.
\item Describe SELinux and AppArmor.
\item Understand container security fundamentals.
\item Identify common operating system vulnerabilities.
\item Answer OS security interview questions confidently.
\end{itemize}

%=========================================================
\section{Security Fundamentals}
%=========================================================

\subsection{Definition}

Security is the protection of computing resources against unauthorized access, modification, destruction, or disclosure.

\subsection{CIA Triad}

The foundation of information security consists of:

\begin{enumerate}
\item Confidentiality
\item Integrity
\item Availability
\end{enumerate}

\begin{table}[h]
\centering
\begin{tabular}{|p{3cm}|p{9cm}|}
\hline
Property & Description \\
\hline
Confidentiality & Prevent unauthorized access to information \\
\hline
Integrity & Prevent unauthorized modification \\
\hline
Availability & Ensure resources remain accessible \\
\hline
\end{tabular}
\caption{CIA Triad}
\end{table}

\subsection{Example}

Online banking systems require:

\begin{itemize}
\item Confidentiality of account data.
\item Integrity of transactions.
\item Availability of services.
\end{itemize}

%=========================================================
\section{Threat Model Basics}
%=========================================================

\subsection{Definition}

A threat model identifies:

\begin{itemize}
\item Assets
\item Threats
\item Attackers
\item Attack Surfaces
\item Mitigations
\end{itemize}

\subsection{Typical Threats}

\begin{itemize}
\item Malware
\item Privilege Escalation
\item Data Theft
\item Insider Threats
\item Remote Exploits
\item Denial of Service
\end{itemize}

\subsection{Security Architecture}

\begin{figure}[h]
\centering
\begin{tikzpicture}

\node[draw] (user) {User};
\node[draw,below=1.5cm of user] (app) {Application};
\node[draw,below=1.5cm of app] (kernel) {Kernel};
\node[draw,below=1.5cm of kernel] (hardware) {Hardware};

\draw[->] (user)--(app);
\draw[->] (app)--(kernel);
\draw[->] (kernel)--(hardware);

\end{tikzpicture}
\caption{Security Layers}
\end{figure}

%=========================================================
\section{Authentication}
%=========================================================

\subsection{Definition}

Authentication verifies identity.

Question answered:

\begin{quote}
Who are you?
\end{quote}

\subsection{Methods}

\begin{itemize}
\item Passwords
\item Biometrics
\item Smart Cards
\item OTPs
\item Multi-Factor Authentication
\end{itemize}

\subsection{Linux Example}

\begin{verbatim}
/etc/passwd
/etc/shadow
\end{verbatim}

\subsection{Windows Example}

Authentication is commonly integrated with:

\begin{itemize}
\item Active Directory
\item Kerberos
\item NTLM
\end{itemize}

%=========================================================
\section{Authorization}
%=========================================================

\subsection{Definition}

Authorization determines allowed actions.

Question answered:

\begin{quote}
What are you allowed to do?
\end{quote}

\subsection{Examples}

\begin{itemize}
\item Read file
\item Write file
\item Execute program
\item Access device
\end{itemize}

\section{Authentication vs Authorization}

\begin{table}[h]
\centering
\begin{tabular}{|p{5cm}|p{5cm}|}
\hline
Authentication & Authorization \\
\hline
Identity Verification & Permission Verification \\
\hline
Who Are You? & What Can You Do? \\
\hline
Occurs First & Occurs After Authentication \\
\hline
Password Login & File Access Check \\
\hline
\end{tabular}
\caption{Authentication vs Authorization}
\end{table}

%=========================================================
\section{Access Control Matrix}
%=========================================================

\subsection{Definition}

An Access Control Matrix defines permissions between subjects and objects.

\begin{table}[h]
\centering
\begin{tabular}{|c|c|c|}
\hline
 & File A & File B \\
\hline
Alice & RW & R \\
\hline
Bob & R & RWX \\
\hline
\end{tabular}
\caption{Access Control Matrix}
\end{table}

\subsection{Challenge}

Matrices become very large in real systems.

%=========================================================
\section{Access Control Lists (ACLs)}
%=========================================================

\subsection{Definition}

ACLs store permissions with each object.

\subsection{Example}

File:

\begin{verbatim}
report.txt
\end{verbatim}

ACL:

\begin{itemize}
\item Alice : Read
\item Bob : Read, Write
\item Charlie : Read
\end{itemize}

\subsection{Advantages}

\begin{itemize}
\item Easy per-object management
\item Widely used
\end{itemize}

\subsection{Disadvantages}

\begin{itemize}
\item Large ACL maintenance
\end{itemize}

%=========================================================
\section{Capabilities}
%=========================================================

\subsection{Definition}

Capabilities associate permissions with subjects instead of objects.

\subsection{Example}

User Alice possesses:

\begin{itemize}
\item Read(File A)
\item Write(File B)
\end{itemize}

\subsection{Advantages}

\begin{itemize}
\item Fine-grained security
\item Easier delegation
\end{itemize}

\subsection{Disadvantages}

\begin{itemize}
\item Capability management complexity
\end{itemize}

\section{ACLs vs Capabilities}

\begin{table}[h]
\centering
\begin{tabular}{|p{5cm}|p{4cm}|p{4cm}|}
\hline
Feature & ACL & Capability \\
\hline
Stored With & Object & Subject \\
\hline
Lookup & Object Based & User Based \\
\hline
Delegation & Difficult & Easier \\
\hline
Usage & Common & Specialized Systems \\
\hline
\end{tabular}
\caption{ACL vs Capability}
\end{table}

%=========================================================
\section{Privilege Separation}
%=========================================================

\subsection{Definition}

Separate privileged and unprivileged components.

\subsection{Example}

Web Browser:

\begin{itemize}
\item Renderer Process
\item Network Process
\item Sandbox Process
\end{itemize}

\subsection{Benefits}

\begin{itemize}
\item Reduced attack impact
\item Better isolation
\end{itemize}

%=========================================================
\section{Principle of Least Privilege}
%=========================================================

\subsection{Definition}

Grant only the minimum permissions required.

\subsection{Example}

A web server should not run as:

\begin{verbatim}
root
\end{verbatim}

Instead:

\begin{verbatim}
www-data
\end{verbatim}

\subsection{Benefits}

\begin{itemize}
\item Reduced damage
\item Improved security
\end{itemize}

%=========================================================
\section{Secure Kernel Design}
%=========================================================

Secure kernels focus on:

\begin{itemize}
\item Isolation
\item Access Control
\item Memory Protection
\item Auditing
\item Secure IPC
\end{itemize}

\subsection{Goals}

\begin{itemize}
\item Small trusted computing base
\item Minimal attack surface
\item Strong privilege boundaries
\end{itemize}

%=========================================================
\section{Memory Protection}
%=========================================================

\subsection{Definition}

Memory protection prevents unauthorized access between processes.

\subsection{Mechanisms}

\begin{itemize}
\item Virtual Memory
\item MMU
\item Page Tables
\item User Mode
\item Kernel Mode
\end{itemize}

\subsection{Benefits}

\begin{itemize}
\item Process Isolation
\item Fault Containment
\item Security Enforcement
\end{itemize}

%=========================================================
\section{ASLR}
%=========================================================

\subsection{Definition}

ASLR stands for Address Space Layout Randomization.

\subsection{Goal}

Randomize memory locations.

\subsection{Protected Areas}

\begin{itemize}
\item Stack
\item Heap
\item Libraries
\item Executable Regions
\end{itemize}

\subsection{Benefit}

Makes exploitation significantly harder.

\subsection{Linux Example}

\begin{verbatim}
cat /proc/sys/kernel/randomize_va_space
\end{verbatim}

%=========================================================
\section{DEP / NX Bit}
%=========================================================

\subsection{Definition}

DEP (Data Execution Prevention) prevents execution from data pages.

NX = No Execute.

\subsection{Purpose}

Stop code execution in:

\begin{itemize}
\item Stack
\item Heap
\item Data Regions
\end{itemize}

\subsection{Mitigates}

\begin{itemize}
\item Buffer Overflow Attacks
\item Code Injection
\end{itemize}

%=========================================================
\section{Sandboxing}
%=========================================================

\subsection{Definition}

Sandboxing isolates applications from the rest of the system.

\subsection{Examples}

\begin{itemize}
\item Chrome Sandbox
\item Android App Sandbox
\item iOS App Sandbox
\item Firejail
\end{itemize}

\subsection{Architecture}

\begin{figure}[h]
\centering
\begin{tikzpicture}

\node[draw] (app) {Application};
\node[draw,below=1.5cm of app] (sandbox) {Sandbox};
\node[draw,below=1.5cm of sandbox] (os) {Operating System};

\draw[->] (app)--(sandbox);
\draw[->] (sandbox)--(os);

\end{tikzpicture}
\caption{Application Sandbox}
\end{figure}

%=========================================================
\section{Secure Boot}
%=========================================================

\subsection{Definition}

Secure Boot verifies cryptographic signatures during boot.

\subsection{Goal}

Prevent unauthorized software from loading.

\subsection{Boot Chain}

\begin{enumerate}
\item Firmware
\item Bootloader
\item Kernel
\item Operating System
\end{enumerate}

\subsection{Benefits}

\begin{itemize}
\item Rootkit Protection
\item Boot-Time Integrity
\end{itemize}

%=========================================================
\section{SELinux Overview}
%=========================================================

\subsection{Definition}

SELinux stands for Security Enhanced Linux.

\subsection{Security Model}

Mandatory Access Control (MAC)

\subsection{Features}

\begin{itemize}
\item Fine-grained policies
\item Process isolation
\item File labeling
\item Security contexts
\end{itemize}

\subsection{Example}

\begin{verbatim}
getenforce
\end{verbatim}

Possible output:

\begin{verbatim}
Enforcing
\end{verbatim}

%=========================================================
\section{AppArmor Overview}
%=========================================================

\subsection{Definition}

AppArmor is a Linux security framework.

\subsection{Approach}

Path-based access control.

\subsection{Advantages}

\begin{itemize}
\item Easier policy management
\item Simpler configuration
\end{itemize}

\subsection{Common Usage}

\begin{itemize}
\item Ubuntu
\item SUSE Linux
\end{itemize}

%=========================================================
\section{Container Security}
%=========================================================

\subsection{Challenges}

Containers share the host kernel.

\subsection{Security Mechanisms}

\begin{itemize}
\item Namespaces
\item Cgroups
\item Capabilities
\item Seccomp
\item SELinux
\item AppArmor
\end{itemize}

\subsection{Container Isolation}

\begin{figure}[h]
\centering
\begin{tikzpicture}

\node[draw] (c1) {Container A};
\node[draw,right=3cm of c1] (c2) {Container B};
\node[draw,below=2cm of c1,xshift=1.5cm] (kernel) {Shared Kernel};

\draw (c1)--(kernel);
\draw (c2)--(kernel);

\end{tikzpicture}
\caption{Container Isolation}
\end{figure}

%=========================================================
\section{Common Vulnerabilities}
%=========================================================

\subsection{Buffer Overflow}

Writing beyond memory boundaries.

\subsection{Privilege Escalation}

Obtaining higher permissions.

\subsection{Race Conditions}

Timing-based attacks.

\subsection{Code Injection}

Injecting malicious code.

\subsection{Directory Traversal}

Accessing unintended files.

\subsection{Use-After-Free}

Accessing freed memory.

%=========================================================
\section{Real-World Attack Scenarios}
%=========================================================

\subsection{Scenario 1}

Web server running as root.

Result:

\begin{itemize}
\item Full system compromise.
\end{itemize}

\subsection{Scenario 2}

No ASLR.

Result:

\begin{itemize}
\item Easier exploitation.
\end{itemize}

\subsection{Scenario 3}

Weak container isolation.

Result:

\begin{itemize}
\item Container escape attack.
\end{itemize}

\subsection{Scenario 4}

Misconfigured ACLs.

Result:

\begin{itemize}
\item Data leakage.
\end{itemize}

%=========================================================
\section{How Modern OSes Isolate Applications}
%=========================================================

Modern operating systems isolate applications using:

\begin{itemize}
\item Virtual Memory
\item Process Boundaries
\item User/Kernel Modes
\item Sandboxing
\item Containers
\item SELinux
\item AppArmor
\item Access Control Systems
\end{itemize}

\subsection{Isolation Layers}

\begin{figure}[h]
\centering
\begin{tikzpicture}

\node[draw] (app1) {App 1};
\node[draw,right=3cm of app1] (app2) {App 2};

\node[draw,below=2cm of app1,xshift=1.5cm] (kernel) {Kernel};

\draw (app1)--(kernel);
\draw (app2)--(kernel);

\end{tikzpicture}
\caption{Application Isolation}
\end{figure}

%=========================================================
\section{Security Best Practices}
%=========================================================

\begin{itemize}
\item Follow least privilege.
\item Keep systems updated.
\item Use strong authentication.
\item Enable ASLR.
\item Enable DEP/NX.
\item Use secure boot.
\item Audit permissions regularly.
\item Isolate workloads.
\item Minimize attack surface.
\item Monitor security logs.
\end{itemize}

%=========================================================
\section{Key Takeaways}
%=========================================================

\begin{itemize}
\item Authentication verifies identity.
\item Authorization verifies permissions.
\item ACLs and capabilities implement access control.
\item Least privilege is a fundamental security principle.
\item Memory protection isolates processes.
\item ASLR randomizes memory layouts.
\item DEP prevents execution from data pages.
\item Sandboxing limits application damage.
\item SELinux and AppArmor strengthen Linux security.
\item Containers require additional security controls.
\end{itemize}

%=========================================================
\section{20 Interview Questions with Answers}
%=========================================================

\begin{enumerate}
\item What is OS security?\\
Answer: Protection of system resources from unauthorized access.

\item What is authentication?\\
Answer: Identity verification.

\item What is authorization?\\
Answer: Permission verification.

\item Difference between authentication and authorization?\\
Answer: Authentication identifies users; authorization determines permissions.

\item What is an ACL?\\
Answer: Object-based permission list.

\item What is a capability?\\
Answer: Subject-held permission token.

\item What is least privilege?\\
Answer: Grant minimum necessary permissions.

\item What is ASLR?\\
Answer: Address Space Layout Randomization.

\item Why use ASLR?\\
Answer: Prevent predictable memory layouts.

\item What is DEP?\\
Answer: Data Execution Prevention.

\item What is NX Bit?\\
Answer: Marks memory non-executable.

\item What is sandboxing?\\
Answer: Application isolation mechanism.

\item What is secure boot?\\
Answer: Verified boot process.

\item What is SELinux?\\
Answer: Mandatory access control framework.

\item What is AppArmor?\\
Answer: Path-based Linux security framework.

\item What is privilege escalation?\\
Answer: Gaining unauthorized privileges.

\item What is container isolation?\\
Answer: Separation using namespaces and cgroups.

\item What is a threat model?\\
Answer: Structured analysis of threats and defenses.

\item Why is memory protection important?\\
Answer: Prevents process interference.

\item What is the CIA triad?\\
Answer: Confidentiality, Integrity, Availability.
\end{enumerate}

%=========================================================
\section{20 MCQs with Answers}
%=========================================================

\begin{enumerate}
\item Authentication verifies:
Answer: Identity

\item Authorization verifies:
Answer: Permissions

\item CIA stands for:
Answer: Confidentiality, Integrity, Availability

\item ACLs are associated with:
Answer: Objects

\item Capabilities are associated with:
Answer: Subjects

\item ASLR randomizes:
Answer: Memory Layout

\item DEP prevents:
Answer: Code Execution in Data Pages

\item NX stands for:
Answer: No Execute

\item SELinux implements:
Answer: Mandatory Access Control

\item AppArmor is:
Answer: Path-Based Security

\item Secure Boot protects:
Answer: Boot Process

\item Least privilege means:
Answer: Minimum Permissions

\item Container isolation uses:
Answer: Namespaces

\item Rootkits commonly target:
Answer: Boot and Kernel Layers

\item Sandboxing provides:
Answer: Isolation

\item Buffer overflow affects:
Answer: Memory Safety

\item Privilege escalation increases:
Answer: Access Rights

\item Threat modeling identifies:
Answer: Risks

\item Virtual memory improves:
Answer: Isolation

\item Secure kernels reduce:
Answer: Attack Surface
\end{enumerate}

%=========================================================
\section{Revision Notes}
%=========================================================

\begin{itemize}
\item Authentication = Identity.
\item Authorization = Permissions.
\item ACL = Object-based permissions.
\item Capability = Subject-based permissions.
\item Least privilege minimizes risk.
\item ASLR randomizes memory addresses.
\item DEP blocks execution from data memory.
\item Sandboxing isolates applications.
\item Secure Boot protects startup integrity.
\item SELinux uses MAC policies.
\item AppArmor uses path-based controls.
\item Containers require layered security.
\end{itemize}

%=========================================================
\section{Cheat Sheet}
%=========================================================

\begin{center}
\begin{tabular}{|p{5cm}|p{8cm}|}
\hline
Authentication & Identity Verification \\
\hline
Authorization & Permission Verification \\
\hline
ACL & Object-Based Permissions \\
\hline
Capability & Subject-Based Permissions \\
\hline
Least Privilege & Minimum Required Access \\
\hline
ASLR & Randomized Memory Layout \\
\hline
DEP & Data Execution Prevention \\
\hline
NX Bit & Non-Executable Memory \\
\hline
Sandbox & Application Isolation \\
\hline
Secure Boot & Verified Startup Chain \\
\hline
SELinux & Mandatory Access Control \\
\hline
AppArmor & Path-Based Access Control \\
\hline
Container Security & Namespaces + Cgroups \\
\hline
Memory Protection & Process Isolation \\
\hline
CIA & Confidentiality Integrity Availability \\
\hline
Threat Model & Assets Threats Defenses \\
\hline
Privilege Separation & Isolated Components \\
\hline
Buffer Overflow & Memory Corruption Attack \\
\hline
Privilege Escalation & Unauthorized Permission Gain \\
\hline
Secure Kernel & Protected Core OS Component \\
\hline
\end{tabular}
\end{center}